\section{The Born--Oppenheimer expansion}\label{app:D.13}

A silicon nucleus is $5.12\times10^4$ times heavier than an electron, with $M = \SI{28.0855}{u} = \SI{26.16}{GeV}$ for the natural isotope mixture, as in \cref{sec:5.3,app:C.6}. Held by force constants of the same electronic order, it vibrates $(M/m_e)^{1/2} \approx 226$ times more slowly than the electrons move, and the electrons have time to settle into their ground state for every position of the nuclei. Born and Oppenheimer turned this into an expansion in powers of $\kappa_{\mathrm{BO}} = (m_e/M)^{1/4}$~\cite{born1927quantentheorie}, in the form that Born and Huang gave it~\cite{born1954dynamical}. Here we derive the equations that \cref{sec:5.1,sec:6.1} state, count every term in powers of $\kappa_{\mathrm{BO}}$, and test the counting on a model that can be solved exactly.

\subsection*{The coupled equations}

We use the Hamiltonian \cref{eq:5.1-H} and the adiabatic states \cref{eq:5.1-adiabatic} of \cref{sec:5.1}, orthonormal at every $u$ and chosen real, and assume that the levels do not cross in the range of $u$ the nuclei explore. The adiabatic states are the states of the clamped-nuclei problem of \cref{chap:3,chap:4}, and the ground-state energy $E_0(u)$ is the energy surface $E^{\mathrm{ad}}$. Any state of electrons and nuclei can be expanded as $\Psi(r,u) = \sum_i\Psi^{\mathrm{nuc}}_i(u)\,\Psi_i(r;u)$, since the adiabatic states are complete in $r$ at each $u$. The nuclear Laplacian acts on both factors, $\nabla^2(\Psi^{\mathrm{nuc}}_i\Psi_i) = (\nabla^2\Psi^{\mathrm{nuc}}_i)\Psi_i + 2\nabla\Psi^{\mathrm{nuc}}_i\cdot\nabla\Psi_i + \Psi^{\mathrm{nuc}}_i\nabla^2\Psi_i$, while $H_{\mathrm{el}}$ acting on $\Psi_i$ yields $E_i(u)\Psi_i$. We insert the expansion into $H\Psi = E\Psi$, multiply by $\Psi_j$ and integrate over the electron coordinates at fixed $u$. By orthonormality the first term of the Laplacian becomes $\delta_{ji}\nabla^2\Psi^{\mathrm{nuc}}_i$; the other two leave brackets that depend on $u$:
\begin{equation*}
  \int\dd r\;\Psi_j\,\nabla^2_{\vect{R}\kappa}\bigl(\Psi^{\mathrm{nuc}}_i\Psi_i\bigr)
  = \delta_{ji}\,\nabla^2_{\vect{R}\kappa}\Psi^{\mathrm{nuc}}_i + 2\,\langle\Psi_j|\nabla_{\vect{R}\kappa}\Psi_i\rangle\cdot\nabla_{\vect{R}\kappa}\Psi^{\mathrm{nuc}}_i + \langle\Psi_j|\nabla^2_{\vect{R}\kappa}\Psi_i\rangle\,\Psi^{\mathrm{nuc}}_i .
\end{equation*}
Multiplied by $-1/2M_\kappa$, summed over the nuclei and over $i$, and added to $E_j(u)\Psi^{\mathrm{nuc}}_j$ from $H_{\mathrm{el}}$, this is the set of coupled equations \cref{eq:5.1-coupled}, with the nonadiabatic operators $\mathcal{N}_{ji}$ of \cref{eq:5.1-N}. The bracket with one gradient is an operator that acts on the nuclear wavefunction through $\nabla_{\vect{R}\kappa}$; the bracket with the Laplacian is a function of $u$ that multiplies it. They are the only terms that connect one electronic state to another, and setting them all to zero is the Born--Oppenheimer approximation \cref{eq:5.1-BO}.

The adiabatic approximation of Born and Huang keeps the diagonal term $\mathcal{N}_{00}$ in addition~\cite{born1954dynamical}. For real states the derivative coupling has no diagonal element, since $\langle\Psi_0|\nabla\Psi_0\rangle = \tfrac12\nabla\langle\Psi_0|\Psi_0\rangle = 0$. Differentiating this once more yields $\langle\Psi_0|\nabla^2\Psi_0\rangle = -\langle\nabla\Psi_0|\nabla\Psi_0\rangle$, so $\mathcal{N}_{00} = \sum_{\vect{R}\kappa}\langle\nabla_{\vect{R}\kappa}\Psi_0|\nabla_{\vect{R}\kappa}\Psi_0\rangle/2M_\kappa$, a small positive potential added to $E^{\mathrm{ad}}$.

\subsection*{The force and the derivative couplings}

Both the force on a nucleus and the derivative couplings follow from one identity. Differentiate \cref{eq:5.1-adiabatic} with respect to $\vect{u}_{\vect{R}\kappa}$ and project on $\Psi_j$:
\begin{equation*}
  \langle\Psi_j|\nabla_{\vect{R}\kappa}H_{\mathrm{el}}|\Psi_i\rangle + E_j\,\langle\Psi_j|\nabla_{\vect{R}\kappa}\Psi_i\rangle = \delta_{ij}\,\nabla_{\vect{R}\kappa}E_i + E_i\,\langle\Psi_j|\nabla_{\vect{R}\kappa}\Psi_i\rangle ,
\end{equation*}
where $H_{\mathrm{el}}$ acting to the left on $\Psi_j$ gave $E_j$. For $j = i$ the derivative couplings cancel, and with $i = 0$ what is left is the Hellmann--Feynman theorem \cref{eq:5.1-HF}, $-\nabla_{\vect{R}\kappa}E^{\mathrm{ad}} = -\langle\Psi_0|\nabla_{\vect{R}\kappa}(V_{eN} + V_{NN})|\Psi_0\rangle$, since only $V_{eN}$ and $V_{NN}$ depend on $u$. For $j \ne i$ the term $\delta_{ij}$ is absent, and $V_{NN}$ drops out because it is a number at fixed $u$ and $\langle\Psi_j|\Psi_i\rangle = 0$. Solving for the bracket yields \cref{eq:5.1-dcoup}: the derivative coupling is the matrix element of the force that the nucleus exerts on the electrons, divided by an electronic excitation energy $E_i - E_j$. Its numerator is the bare electron--phonon vertex, a one-body operator, $\nabla_{\vect{R}\kappa}V_{eN} = \sum_e\nabla_{\vect{R}\kappa}v_\kappa(\vect{r}_e - \vect{R} - \vect{\tau}_\kappa - \vect{u}_{\vect{R}\kappa})$ with $v_\kappa$ the electron--nucleus potential. Its denominator decides when the adiabatic picture fails: the coupling is small when the excitation energies are large, and it diverges where two adiabatic levels meet.

\subsection*{Power counting}

We measure lengths in $a_0$ and energies in $\alpha^2m_e$, as in \cref{sec:5.1}. The balance \cref{eq:5.1-scales} then makes the zero-point amplitude the fraction $\kappa_{\mathrm{BO}}$ of the Bohr radius, $\ell_{\mathrm{zp}} \sim \kappa_{\mathrm{BO}}a_0$, and a vibrational quantum the fraction $\kappa_{\mathrm{BO}}^2$ of an electronic energy; \cref{we:5-kappa} gives the numbers for silicon. If the displacement is written in units of $\kappa_{\mathrm{BO}}a_0$, the nuclear kinetic energy and the harmonic potential are both $\sim\kappa_{\mathrm{BO}}^2\alpha^2m_e$, and every other term is smaller by further powers of $\kappa_{\mathrm{BO}}$.

The remaining terms follow by the substitution $u \sim \kappa_{\mathrm{BO}}a_0$ and $\nabla \sim 1/(\kappa_{\mathrm{BO}}a_0)$ acting on a nuclear wavefunction, with force constants $\Phi^{(n)} \sim \alpha^2m_e/a_0^n$ and $1/M = \kappa_{\mathrm{BO}}^4/m_e$. The cubic term of the energy surface is $\sim\alpha^2m_e\,\kappa_{\mathrm{BO}}^3$, smaller than a phonon energy by one power of $\kappa_{\mathrm{BO}}$, which is the anharmonic small parameter of \cref{sec:5.7}. The quartic term is $\sim\alpha^2m_e\,\kappa_{\mathrm{BO}}^4$, and so is the cubic term at second order, $(\kappa_{\mathrm{BO}}^3)^2/\kappa_{\mathrm{BO}}^2$, since its energy denominator is a phonon energy: the anharmonic shifts of the vibrational levels are $\sim\alpha^2m_e\,\kappa_{\mathrm{BO}}^4$. In the nonadiabatic operators \cref{eq:5.1-N}, the electronic brackets vary on the scale $a_0$, $\langle\Psi_j|\nabla\Psi_i\rangle \sim 1/a_0$ and $\langle\Psi_j|\nabla^2\Psi_i\rangle \sim 1/a_0^2$. The first-derivative coupling is then $(1/M)(1/a_0)(1/\kappa_{\mathrm{BO}}a_0) = \alpha^2m_e\,\kappa_{\mathrm{BO}}^3$, and the second-derivative coupling $\alpha^2m_e\,\kappa_{\mathrm{BO}}^4$. The diagonal correction $\mathcal{N}_{00}$ is $\sim\kappa_{\mathrm{BO}}^4$ as well, but most of it is a constant. Its variation over a zero-point amplitude changes the curvature of the surface by the fraction $\kappa_{\mathrm{BO}}^4$, and the vibrational levels by $\alpha^2m_e\,\kappa_{\mathrm{BO}}^6$. The derivative coupling connects the electronic ground state to excited states at an energy $\sim\alpha^2m_e$ above it, and it therefore shifts the vibrational levels only at second order, by $(\kappa_{\mathrm{BO}}^3)^2\alpha^2m_e = \alpha^2m_e\,\kappa_{\mathrm{BO}}^6$. Every nonadiabatic correction to the vibrational energies is therefore of relative order $\kappa_{\mathrm{BO}}^4 = m_e/M$, two powers of $\kappa_{\mathrm{BO}}$ beyond the anharmonic ones. The electronic energy, the vibrations and the anharmonic corrections appear at orders $\kappa_{\mathrm{BO}}^0$, $\kappa_{\mathrm{BO}}^2$ and $\kappa_{\mathrm{BO}}^4$, and the coupling of the electronic and nuclear motion first at $\kappa_{\mathrm{BO}}^6$~\cite{born1927quantentheorie}.

The counting assumed that the electronic excitation energies $\Delta E$ are of order $\alpha^2m_e$. When they are smaller, the derivative coupling \cref{eq:5.1-dcoup} grows as $1/\Delta E$. Its matrix element between vibronic states is $(1/M)\langle\Psi_j|\nabla\Psi_0\rangle\cdot\nabla \sim (1/M)(F_{j0}/\Delta E)(1/\ell_{\mathrm{zp}})$, with $F_{j0}$ the matrix element of the force in the numerator of \cref{eq:5.1-dcoup}; with $1/M = \omega_{\mathrm{ph}}\ell_{\mathrm{zp}}^2$ for $\ell_{\mathrm{zp}}^2 = 1/(M\omega_{\mathrm{ph}})$ it is $\omega_{\mathrm{ph}}\,\ell_{\mathrm{zp}}F_{j0}/\Delta E$. Divided by the energy $\Delta E$ that separates the two vibronic states, it gives the amplitude with which an excited electronic state is mixed into the ground state,
\begin{equation}
  \frac{\omega_{\mathrm{ph}}}{\Delta E}\;\frac{\ell_{\mathrm{zp}} F_{j0}}{\Delta E} ,
  \label{eq:D.13-mixing}
\end{equation}
which is \cref{eq:5.1-mixing}. At $\Delta E \sim \alpha^2m_e$ the two factors are $\kappa_{\mathrm{BO}}^2$ and $\kappa_{\mathrm{BO}}$, the $\kappa_{\mathrm{BO}}^3$ above.

\subsection*{Clamped-nuclei states}

The adiabatic states follow the nuclei, and the price is the derivative couplings. The alternative is to expand in the electronic states at the equilibrium configuration, $\Psi(r,u) = \sum_i\tilde\Psi^{\mathrm{nuc}}_i(u)\,\Psi_i(r;0)$, which do not depend on $u$ and have no derivative couplings. The coupling then sits in the Hamiltonian, $H_{\mathrm{el}}(u) = H_{\mathrm{el}}(0) + [V_{eN}(u) - V_{eN}(0)] + [V_{NN}(u) - V_{NN}(0)]$, and its expansion in the displacements is
\begin{equation}
  V_{eN}(u) - V_{eN}(0) = \sum_{\vect{R}\kappa}\vect{u}_{\vect{R}\kappa}\cdot\nabla_{\vect{R}\kappa}V_{eN}\Big|_0
  + \frac12\sum_{\vect{R}\kappa,\vect{R}'\kappa'}\vect{u}_{\vect{R}\kappa}\cdot\bigl(\nabla_{\vect{R}\kappa}\nabla_{\vect{R}'\kappa'}V_{eN}\bigr)\Big|_0\cdot\vect{u}_{\vect{R}'\kappa'} + \cdots .
  \label{eq:D.13-expand}
\end{equation}
The linear term is a one-body electron operator times a displacement. Written in second quantization over the Bloch states of the clamped crystal (\cref{eq:2.3-onebody}), with $\vect{u}_{\vect{R}\kappa}$ expanded in phonons by \cref{eq:A.6-u}, it has the form \cref{eq:A.6-Hep} of the electron--phonon Hamiltonian. Its vertex is \cref{eq:A.6-g} with the bare derivative $\nabla V_{eN}$ in place of $\partial_{\kappa\alpha,\vect{Q}}V_{\mathrm{KS}}$. The other electrons screen this bare vertex, and in density-functional theory the self-consistent change of the Kohn--Sham potential replaces it (\cref{sec:6.2}). The quadratic term is the second-order (Allen--Heine) coupling of \cref{sec:6.3}. The clamped-nuclei approximation of \cref{chap:3,chap:4} keeps $H_{\mathrm{el}}(0)$ alone.

The two expansions are equivalent order by order in $u$. At first order, the linear term of \cref{eq:D.13-expand} mixes into $\Psi_i(r;0)$ the states $\Psi_j(r;0)$ with amplitude $\vect{u}\cdot\langle\Psi_j|\nabla V_{eN}|\Psi_i\rangle/(E_i - E_j)$, and by \cref{eq:5.1-dcoup} this is $\vect{u}\cdot\langle\Psi_j|\nabla\Psi_i\rangle$, the first term of the Taylor series of $\Psi_i(r;u)$. The adiabatic states are the clamped states dressed by the electron--phonon coupling to all orders in $u$, with the nuclei held still; the nonadiabatic operators are what remains of the coupling when the nuclei move.

\subsection*{A solvable harmonic model}

One electron bound harmonically to one nucleus, which is itself bound to a fixed point, tests the counting without approximation. In one dimension,
\begin{equation*}
  H = \frac{p^2}{2m_e} + \frac{P^2}{2M} + \frac{k_{\mathrm{e}}}{2}(x - X)^2 + \frac{K_{\mathrm{N}}}{2}X^2 ,
\end{equation*}
with $x$ and $p$ the electron's coordinate and momentum, $X$ and $P$ the nucleus's, and two spring constants of the same electronic order. The adiabatic states are oscillator states centered on the nucleus, with the excitation energy $\omega_{\mathrm{e}} = (k_{\mathrm{e}}/m_e)^{1/2}$ and the surface $E_0(X) = \omega_{\mathrm{e}}/2 + K_{\mathrm{N}}X^2/2$, so the Born--Oppenheimer frequency is $\omega_{\mathrm{BO}} = (K_{\mathrm{N}}/M)^{1/2}$. The diagonal correction $\mathcal{N}_{00} = m_e\omega_{\mathrm{e}}/(4M)$ is a constant, since the ground state only translates with $X$, and it leaves the frequency unchanged. The exact frequencies are those of two coupled oscillators. With the force constants as a matrix, the equation of motion requires $\det\bigl[\begin{smallmatrix}k_{\mathrm{e}} - m_e\omega^2 & -k_{\mathrm{e}}\\ -k_{\mathrm{e}} & k_{\mathrm{e}} + K_{\mathrm{N}} - M\omega^2\end{smallmatrix}\bigr] = 0$, that is $m_eM\omega^4 - [m_e(k_{\mathrm{e}} + K_{\mathrm{N}}) + Mk_{\mathrm{e}}]\,\omega^2 + k_{\mathrm{e}}K_{\mathrm{N}} = 0$. Its lower root, expanded in $m_e/M$, is
\begin{equation}
  \omega^2 = \frac{K_{\mathrm{N}}}{M}\Bigl[1 - \frac{m_e}{M} + \Bigl(1 - \frac{K_{\mathrm{N}}}{k_{\mathrm{e}}}\Bigr)\frac{m_e^2}{M^2} + O\Bigl(\frac{m_e^3}{M^3}\Bigr)\Bigr] .
  \label{eq:D.13-model}
\end{equation}
At first order the nucleus vibrates as if its mass were $M + m_e$: the electron is dragged along and adds its mass. Its lag behind the nucleus, which depends on the ratio of the two springs, enters only at the next order. The relative correction is $m_e/M = \kappa_{\mathrm{BO}}^4$ for $\omega^2$, and $m_e/2M$ for the frequency itself, and it comes entirely from the off-diagonal coupling at second order, as the counting predicts. For $K_{\mathrm{N}} = k_{\mathrm{e}}/2$ and $m_e/M = 10^{-2}$, $10^{-3}$ and $10^{-4}$, the exact $\omega^2$ differs from $(K_{\mathrm{N}}/M)(1 - m_e/M)$ by $5.1\times10^{-5}$, $5.0\times10^{-7}$ and $5.0\times10^{-9}$ of its value, the $m_e^2/M^2$ term of \cref{eq:D.13-model} plus, at the largest mass ratio, a visible $m_e^3/M^3$ term.
